
% !TEX TS-program = lualatexmk
% !TEX parameter =  --shell-escape

\documentclass[vespers_book_la_eng.tex]{subfiles}

\ifcsname preamble@file\endcsname
  \setcounter{page}{\getpagerefnumber{M-vespers_book_la_eng_trinity_sc}}
\fi

\begin{document}

\litdateen{I Sunday after Pentecost}
\festum{feast of the most holy trinity}
\addcontentsline{toc}{chapter}{Most Holy Trinity through the octave of the Sacred Heart}\phantomsection
\header{feast of the trinity through the octave of the sacred heart.}

\doubleclass{I}

\atvespers{i}

\lastplace{II}{116}

\psalmus{116}{116_5a}{116_5}{116_en}

\label{cap_trin}
\chapterlateng
\scripture{Rom. 11, 33.}

\texteparacol{\lettrine{O}{} Altitúdo divitiárum sapiéntiæ et sciéntiæ Dei:~† quam incomprehensibília sunt judícia ejus,~* et investigábiles viæ ejus!}{english}
{\noindent O the depth of the riches of the wisdom and of the knowledge of God! How incomprehensible are his judgments, and how unsearchable his ways!}

\label{hymn_trin}
\hymnlateng

\gscore[]{8.}{hy_jam_sol_recedit_igneus_blessed_trinity_solesmes_1961}
\textes{jam_sol_trinity_en}

\smallscore{versiculus_1_vesperis_trinity}
\begin{enpars}\noindent \vv Let us bless the Father, and the Son with the Holy Ghost.
 
\noindent \rr Let us praise him and exalt him above all forever.\end{enpars}

\gscore[Ad Magnif.]{Ant. 1. D}{an_gratias_tibi_deus_solesmes_1961}
\antiphona{english}{Thanks be to thee, o God, thanks be unto thee, o true and one Trinity, o one almighty deity, holy and one Unity}

\label{oratio_trin}\phantomsection

\collectlateng

\texteparacol{\lettrine{O}{m}nípotens sempitérne Deus, qui dedísti fámulis tuis in confessióne veræ fídei, ætérnæ Trinitátis glóriam agnóscere, et in poténtia majestátis adoráre Unitátem:~† quǽsumus; ut ejúsdem fídei firmitáte,~* ab ómnibus semper muniámur advérsis. Per Dóminum.}{english}
{\noindent
Almighty and everlasting God, thou hast given thy servants in the confession of the true faith, to acknowledge the glory of the eternal Trinity,  and to adore the unity in the power of thy majesty; grant, that by steadfastness in this faith we may ever be defended from all adversities.
Through our Lord.}

\rubrique{A commemoration of the I Sunday after Pentecost is always made.}

\gscore[Ant.]{1.}{an_loquere_domine_solesmes_1961}
\antiphona{english}{Speak, Lord, for thy servant heareth}
 
        \texteparacol{\vv Vespertína orátio ascéndat ad te Dómine.

\rr Et descéndat super nos misericórdia tua.}{english}{
 \noindent \vv Let the evening prayer ascend unto thee, o Lord.
                  
\noindent \rr And let there descend upon us thy mercy.}
\label{oratio_1_dom_pp}
\collectlateng

  \texteparacol{\lettrine{D}{e}us, in te sperántium fortitúdo, adésto propítius invocatiónibus nostris:~† et quia sine te nihil potest mortális infírmitas, præsta auxílium grátiæ tuæ;~* ut in exsequéndis mandátis tuis, et voluntáte tibi et actióne placeámus. Per Dóminum.}
{english}
{\noindent
O God, the Strength of all them that put their trust in thee, mercifully accept our prayers, and because through the weakness of our mortal nature we can do no good thing without thee, grant us the help of thy grace, that in keeping of thy commandments we may please thee, both in will and deed.
Through our Lord.}

\invesperis{ii}

\firstantiphon{1. f}

\initialscore{an_gloria_tibi_trinitas_solesmes_1961}
\antiphona{english}{Glory be to thee, o equal Trinity one deity, from before all ages, so now and for evermore}
\psalmus{109}{109_1f}{109_1}{109_en}

 \gscore[2. Ant.]{2. D}{an_laus_et_perennis_gloria_solesmes_1961}
\antiphona{english}{Praise and endless glory be to God the Father, the Son, and the Holy Paraclete, for ever and ever}
\psalmus{110}{110_2}{110_2}{110_en}

  \gscore[3. Ant.]{3. a2}{an_gloria_laudis_resonet_solesmes_1961}
 \antiphona{english}{Let the glory of praise be heard from the mouth of all to the Father, and to the begotten Son; let it equally resound with perpetual praise to the Holy Spirit}
 \psalmus{111}{111_3a2}{111_3a2_g}{111_en}

\gscore[4. Ant.]{4. E}{an_laus_deo_patri_solesmes_1961}
 \antiphona{english}{Praise be to God the Father, and to his co-equal Son, and to thee, O Holy Ghost: with perpetual delight may it resound from our lips for ever}
\psalmus{112}{112_4E}{112_4E}{112_en}

 \gscore[5. Ant.]{5. a}{an_ex_quo_omnia_solesmes_1961}
 \antiphona{english}{From whom are all things, through whom are all things, and in whom are all things: to him be glory for evermore}
\psalmus{113}{113_5a}{113_5}{113_en}

\caphymnen{I}

\smallscore{versiculus_2_vesperis_trinity}

  \begin{enpars}\noindent \vv Blessed art thou, O Lord, in the firmament of heaven.
   
\noindent \rr And worthy of praise, and glorious for ever.\end{enpars}

 \gscore[Ad Magnif.]{Ant. 4. E}{an_te_deum_patrem_solesmes}
\antiphona{english}{With all our heart and with all our voice do we acknowledge thee, praise thee, and bless thee, o God the Father the Unbegotten, o God the Only-Begotten Son, o God the Holy Ghost the Comforter, one holy and undivided Trinity, to thee be glory forever}
 
 \rubrique{Collect \normaltext{Omnípotens sempitérne Deus, \pageref{oratio_trin}} as at I Vespers.}
 
 \rubrique{For the commemoration of the I Sunday after Pentecost:}
 
 \gscore[Ant.]{8.}{an_nolite_judicare_solesmes_1961}
\antiphona{english}{Judge not, that ye be not judged, for with what judgment ye judge, ye shall be judged, saith the Lord}
 
        \texteparacol
{\vv Dirigátur Dómine orátio mea.

\rr Sicut incénsum in conspéctu tuo.}
{english}
{\noindent \vv Direct, O Lord, my prayer.
                  
\noindent \rr As incense in thy sight.}

 \rubrique{Collect \normaltext{Deus, in te sperántium fortitúdo, \pageref{oratio_1_dom_pp}}.}

\litdateen{Thursday of the I Week after the Octave of Pentecost}
\festum{feast of corpus christi}
\addcontentsline{toc}{section}{Corpus Christi} \phantomsection
\label{corpus_christi}
\doubleclass{I}[with a Privileged Octave of the III Order]
 
\atvespers{i}

\firstantiphon{1. f}
\initialscore{an_sacerdos_in_aeternum_solesmes_1961}

\antiphona{english}{Christ the Lord, being made an High Priest for ever after the order of Melchisedec, hath offered bread and wine}

\psalmus{109}{109_1f}{109_1}{109_en}

 \gscore[2. Ant.]{2. D}{an_miserator_dominus_solesmes_1961}
\antiphona{english}{He hath made his wonderful works to be remembered; the Lord is gracious and full of compassion. He hath given meat unto them that fear him}
\psalmus{110}{110_2}{110_2}{110_en}

  \gscore[3. Ant.]{3. a2}{an_calicem_salutaris_accipiam_et_sacrificabo_solesmes_1961}
 \antiphona{english}{I will take the cup of salvation, and offer the sacrifice of thanksgiving}
 \psalmus{115}{115_3a2}{115_3a2_g}{115_en}
 
\gscore[4. Ant.]{4. E}{an_sicut_novellae_solesmes_1961}
\antiphona{english}{Let the children of the Church be like olive-plants round about the table of the Lord}
\psalmus{127}{127_4E}{127_4E}{127_en}

 \gscore[5. Ant.]{5. a}{an_qui_pacem_ponit_solesmes_1961}
\antiphona{english}{The Lord, That maketh peace in the borders of the Church, filleth her with the finest of the wheat}
\psalmus{147}{147_5}{147_5}{147_en}

\chapterlateng
 \scripture{1 Cor. 11, 23 – 24.}
 
\texteparacol{\lettrine{F}{r}atres: Ego enim accépi a Dómino quod et trádidi vobis,~† quóniam Dóminus Jesus, in qua no\-cte tradebátur, accépit panem, et grátias agens fregit, et dixit: Accípite et manducáte: hoc est corpus meum, quod pro vobis tradétur:~* hoc fácite in meam commemoratiónem.}{english}{\noindent Brethren: For I have received of the Lord that which also I delivered unto you, that the Lord Jesus, the same night in which he was betrayed, took bread. And giving thanks, broke, and said: Take ye, and eat: this is my body, which shall be delivered for you: this do for the commemoration of me.}

\hymnlateng  \label{hy_pange_lingua_corpus_christi_solesmes} \phantomsection

\rubrique{If the Office is celebrated before the exposed Blessed Sacrament, the verses beginning at \normaltext{Tantum ergo sacraméntum} are sung kneeling.}

\gscore[]{3.}{hy_pange_lingua_corpus_christi_solesmes}

\textes{pange_lingua_en}

\textes{octave_versiculus_rubrique}

\label{panem}\phantomsection
   \smallscore{versiculus_vesperis_cc}
   
\begin{enpars}
\noindent \vv Thou hast given them bread from heaven, alleluia.

\noindent \rr Having all sweetness within it, alleluia.

\end{enpars}
%\myrule

\label{o_quam_suavis} \phantomsection
\gscore[Ad Magnif.]{6. F}{an_o_quam_suavis_solesmes_1961}
\antiphona{english}{O how sweet is thy spirit, Lord,
thou who, in order to demonstrate thy sweetness to thy children,
send down from heaven the sweetest bread unsurpassed,
filling the hungry with good things,
sending away empty the disdainful rich}

\label{oratio_cc}\phantomsection

\collectlateng
\texteparacol{\lettrine{D}{e}us, qui nobis sub Sacraménto mirábili passiónis tuæ memóriam reliquísti:~† tríbue, quǽsumus, ita nos córporis, et sánguinis tui sacra mystéria venerári;~* ut redem\-ptiónis tuæ fru\-ctum in nobis júgiter sentiámus. Qui vivis et regnas cum Deo Patre in unitáte.}{english}
{\noindent Let us pray. O God, who under a wonderful Sacrament hast left us a memorial of thy Passion; grant us, we beseech thee, so to reverence the sacred mysteries of thy Body and Blood, that we may ever feel within ourselves the fruit of thy redemption. Who livest and reignest with God the Father in the unity.}

\atvespers{ii}

 \allexcept

%\caphymnen{I}

 \label{an_o_sacrum_convivium_solesmes_1961}\phantomsection
 \gscore[Ad Magnif.]{Ant. 5. a}{an_o_sacrum_convivium_solesmes_1961}
 
\antiphona{english}{O Sacred Banquet, in which Christ is received, the memory of his Passion is recalled, the soul is filled with grace, and the pledge of future glory is given unto us}

\rubrique{The Office is celebrated thusly for the entire octave.}

    \bigtitle{Saturday within the Octave of Corpus Christi.}
    \addcontentsline{toc}{section}{Saturday within the octave of Corpus Christi} \phantomsection
    
    \chapterlateng
    \label{cap_dom_infra_oct_cc} \phantomsection
    \scripture{1 Joann. 3, 13 – 14.}

\texteparacol{\lettrine{C}{a}ríssimi: Nolíte mirári, si odit vos mundus.~† Nos scimus quóniam transláti sumus de morte ad vitam,~* quóniam dilígimus fratres.}{english}{\noindent Beloved: Wonder not, brethren, if the world hate you. We know that we have passed from death to life, because we love the brethren. He that loveth not, abideth in death.}

\rubrique{Hymn \normaltext{Pangue lingua, \pageref{hy_pange_lingua_corpus_christi_solesmes}.}}

\label{versiculus_dom_octave_cc} \phantomsection
 \smallscore{versiculus_dom_octave_cc}

\begin{enpars}\noindent \vv He fed them with the finest of the wheat, alleluia.

\noindent \rr  And with honey out of the rock did He satisfy them, alleluia.
\end{enpars}

        \gscore[Ad Magnif.]{Ant. 7. a}{an_puer_samuel_solesmes_1961}
   \antiphona{english}{The child Samuel ministered unto God before Eli, and the word of the Lord was precious in his sight}
   
   \label{oratio_dom_cc}\phantomsection

\collectlateng

\texteparacol{\lettrine{S}{a}ncti nóminis tui, Dómine, timórem páriter et amórem fac nos habére perpétuum:~† quia nunquam tua gubernatióne destítuis,~* quos in soliditáte tuæ dilectiónis instítuis.
Per Dóminum.}{english}{
\noindent Let us pray.
O Lord, who never failest to help and govern them whom thou dost bring up in thy steadfast fear and love keep us, we beseech thee, under the protection of thy good providence, and make us to have a perpetual fear and love of thy holy Name.
Through our Lord.}
   
   \rubrique{On Saturday, a commemoration is made of the octave, Ant. \normaltext{O sacrum convívium, \pageref{an_o_sacrum_convivium_solesmes_1961}, ℣. Panem de cælo, \pageref{panem}.} Collect \normaltext{Deus qui nobis, \pageref{oratio_cc}.}}
    
    \bigtitle{Sunday within the Octave of Corpus Christi.}
    
    {\centering which is the II after Pentecost.\par}
\addcontentsline{toc}{section}{Sunday within the octave of Corpus Christi} \phantomsection

\rubrique{Antiphons and psalms as at I Vespers of the feast, \normaltext{\pageref{corpus_christi}.}}

\rubrique{Chapter \normaltext{Carissime: Nolíte.} as above, \normaltext{\pageref{cap_dom_infra_oct_cc}.}}

\rubrique{Hymn \normaltext{Pangue lingua, \pageref{versiculus_dom_octave_cc}.}}

\rubrique{\normaltext{℣. Cibávit,} as above, \normaltext{\pageref{hy_pange_lingua_corpus_christi_solesmes}.}}

      \gscore[Ad Magnif.]{Ant. 5. a}{an_exi_cito_in_plateas_solesmes_1961}
       \antiphona{english}{Go out quickly into the streets and lanes of the city, and compel the poor and the maimed, and the halt, and the blind, to come in that my house may be full, alleluia}
       
       \rubrique{\begin{enpars}Then the following day is commemorated, Ant. \normaltext{O quam suávis, \pageref{o_quam_suavis}. ℣. Panem, \pageref{panem},} Collect \normaltext{Deus, qui nobis, \pageref{oratio_cc}.} But if the office is not of the following day within the octave, then the commemoration is made as follows, Ant. \normaltext{O sacrum convívium, \pageref{an_o_sacrum_convivium_solesmes_1961}, ℣. Panem de cælo, \pageref{panem},} Collect \normaltext{Deus, qui nobis, \pageref{oratio_cc}.}\end{enpars}} 

\rubrique{During the octave, the office is as on the feast, unless a double of the I class feast is celebrated, and is of semidouble rite.}

\rubrique{On Wednesday, all is at I Vespers, and the office is of double major rite.}

\rubrique{If however the feast of the Nativity of Saint John the Baptist or that of the Holy Apostles Peter and Paul falls on the following day, then at I Vespers, for the commemoration of the preceding day within the octave is said the antiphon \normaltext{O sacrum convívium,} \pageref{an_o_sacrum_convivium_solesmes_1961} as at II Vespers of the feast.}

\newpage

\litdateen{Friday after the Octave of Corpus Christi}
\festum{feast of the most sacred heart of jesus}
\addcontentsline{toc}{section}{Most Sacred Heart of Jesus} \phantomsection

\doubleclass{I}[with a Privileged Octave of the III Order]
 
\atvespers{i}

\firstantiphon{1. g}
\zerobaroffsettextleft
 \initialscore{an_suavi_jugo_solesmes_1961} %% watch in case initial remains low after a couple runs
 \resetbaroffsettextleft
    \antiphona{english}{Put thy easy yoke upon all mankind, O Lord, and be thou ruler, even in the midst of thine enemies}
\psalmus{109}{109_1g}{109_1}{109_en}
 
 \gscore[2. Ant.]{2. D}{an_misericors_et_miserator_solesmes_1961}
      \antiphona{english}{This merciful and gracious Lord hath given Meat unto them that fear him}
 \psalmus{110}{110_2}{110_2}{110_en}

 
%  \rubrique{Ps. \normaltext{Confitébor tibi, \pageref{M-110_2}.}}
 
  \gscore[3. Ant.]{7. a}{an_exortum_est_sacred_heart_solesmes_1961}
        \antiphona{english}{Unto the godly there ariseth up light in the darkness, even the Lord merciful and gracious}
  \psalmus{111}{111_7a}{111_7}{111_en}

   \gscore[4. Ant.]{8. c}{an_quid_retribuam_solesmes_1961}
         \begin{enpars}\noindent Ant. What reward shall I give unto the Lord for all the benefits that he hath done unto me?\end{enpars}
  \psalmus{115}{115_8c}{115_8}{115_en}
   
    \gscore[5. Ant.]{4. A*}{an_apud_dominum_propitiatio_solesmes_1961}
          \antiphona{english}{With the Lord there is mercy, and with him is plenteous redemption}
    \psalmus{129}{129_4A_A_star}{129_4A_A_star}{129_en}
    
    \chapterlateng

%\smalltitle{On the feast and within the octave:}
 \scripture{Eph. 3, 8 – 9.}

\texteparacol{\lettrine{F}{r}ratres: Mihi ómnium san\-ctórum mínimo data est grátia hæc,~† in géntibus evangelizáre investigábiles divítias Christi:~* et illumináre omnes, quæ sit dispensátio sacraménti abscónditi a sǽculis in Deo.}{english}{\noindent Brethren: Unto me, who am less than the least of all saints, is this grace given, that I should preach among the Gentiles the unsearchable riches of Christ; and to make all men see what is the fellowship of the Mystery, which from the beginning of the world hath been hid in God.}

\label{hy_en_ut_superba_solesmes_1961} \phantomsection
\hymnlateng
\gscore[]{3.}{hy_en_ut_superba_solesmes_1961}
\textes{en_ut_superba_en}

   \smallscore{versiculus_i_vesperis_ss_cordis}\label{tollite}\phantomsection
   \begin{enpars}
\noindent \vv Take my yoke upon you, and learn of me.

\noindent \rr For I am meek and humble of heart.
\end{enpars}
   
\label{Ignem}
\gscore[Ad Magnif,]{1. D}{an_ignem_solesmes_1961}
  \begin{enpars}\noindent Ant. I am come to send Fire on the earth, and what will I, but that it be kindled?\end{enpars}

\label{oratio_sc}\phantomsection

\collectlateng
\texteparacol{\lettrine{D}{e}us, qui nobis in Corde Fílii tui, nostris vulneráto peccátis, infinítos dilectiónis thesáuros misericórditer largíri dignáris:~† concéde quǽsumus; ut illi devótum pietátis nostræ præstántes obséquium,~* dignæ quoque satisfactiónis exhibeámus offícium. Per eúmdem Dóminum.}{english}
{\noindent Let us pray. O God, who hast suffered the heart of thy Son to be wounded by our sins, and in that very heart hast bestowed on us the abundant riches of thy love: grant that the devout homage of our hearts, which we render unto him, may of thy mercy be deemed a recompence acceptable in thy sight.
Through the same Lord.}

\rubrique{No commemorations are made, except only that of the feast of the Nativity of Saint John the Baptist or that of the Holy Apostles Peter and Paul, if it falls on the preceding day.}

\atvespers{ii}
\firstantiphon{1. f}
\initialscore{an_unus_militum_solesmes_1961}
        \antiphona{english}{One of the soldiers with a spear pierced his side, and forthwith came there out Blood and Water}
\psalmus{109}{109_1f}{109_1}{109_en}

 \gscore[2. Ant.]{7. c}{an_stans_jesus_solesmes_1961}
         \antiphona{english}{Jesus stood and cried, saying, If any man thirst, let him come unto me and drink}
\psalmus{110}{110_7c}{110_7}{110_en}

  \gscore[3. Ant.]{3. a2}{an_in_caritate_solesmes_1961}
          \antiphona{english}{With an everlasting love hath God loved us, and from the day that he was lifted up over the earth, he hath drawn us with loving-kindness unto his Heart}
 \psalmus{115}{115_3a2}{115_3a2_g}{115_en}

\gscore[4. Ant.]{4. E}{an_venite_ad_me_solesmes_1961}
        \antiphona{english}{Come unto me, all ye that travail and are heavy laden, and I will refresh you}
\psalmus{127}{127_4E}{127_4E}{127_en}

 \gscore[5. Ant.]{5. a}{an_fili_praebe_solesmes_1961}
        \antiphona{english}{My son, give me thine heart, and let thine eyes observe my ways}
\psalmus{147}{147_5}{147_5}{147_en}

\caphymnen{I}

%    \chapterlateng

%\smalltitle{On the feast and within the octave:}
% \scripture{Eph. 3, 8 – 9.}
%
%\texteparacol{\lettrine{F}{r}ratres: Mihi ómnium san\-ctórum mínimo data est grátia hæc,~† in géntibus evangelizáre investigábiles divítias Christi:~* et illumináre omnes, quæ sit dispensátio sacraménti abscónditi a sǽculis in Deo.}{english}{\noindent Brethren: Unto me, who am less than the least of all saints, is this grace given, that I should preach among the Gentiles the unsearchable riches of Christ; and to make all men see what is the fellowship of the Mystery, which from the beginning of the world hath been hid in God.}

%\rubrique{Hymn \normaltext{En ut supérba críminum, \pageref{hy_en_ut_superba_solesmes_1961}.}}

\textes{octave_versiculus_rubrique}

\label{haurietis}\phantomsection
   \smallscore{versiculus_ii_vesperis_ss_cordis}
   
\begin{enpars}
\noindent \vv With joy shall ye draw water.

\noindent \rr Out of the fountains of the Saviour.
\end{enpars}

 \label{ad_jesum_autem}\phantomsection
 \gscore[]{1. f}{an_ad_jesum_autem_solesmes_1961}
 
 \begin{enpars}{english} \noindent Ant. But when they came to Jesus and saw that he was dead already, they broke not his legs, but one of the soldiers with a spear pierced his side, and forthwith came there out Blood and Water.\end{enpars}

    \bigtitle{Saturday within the octave of the Sacred Heart.}
\addcontentsline{toc}{section}{Saturday within the octave of the Sacred Heart} \phantomsection

    \chapterlateng
\label{cap_dom_infra_oct_sc} \phantomsection
\scripture{1 Petri 5, 6 – 7.}

\texteparacol{\lettrine{C}{a}ríssimi: Humiliámini sub poténti manu Dei, ut vos exáltet in témpore visitatiónis;~† omnem sollicitúdinem vestram proiciéntes in eum,~* quóniam ipsi cura est de vobis.}{english}{\noindent Beloved:  humble yourselves therefore under the mighty hand of God, that he may exalt you in due time: casting all your care upon him; for he careth for you.}

\rubrique{Hymn \normaltext{En ut supérba críminum, \pageref{hy_en_ut_superba_solesmes_1961}.}}

\label{versiculus_dom_octave_sc}\phantomsection

 \smallscore{versiculus_dom_octave_sc}

\begin{enpars}\noindent \vv The merciful Lord hath instituted a memorial of his wondrous deeds.

\noindent \rr  He hath given Meat unto them that fear him.\end{enpars}

\gscore[Ad Magnif.]{Ant. 1. f}{an_cognoverunt_omnes_solesmes_1961}
\antiphona{english}{And they all knew, from Dan even unto Beersheba, that Samuel was established to be a Prophet of the Lord}

\label{oratio_dom_sc}\phantomsection
\collectlateng
\texteparacol{\lettrine{P}{r}oté\-ctor in te sperántium Deus, sine quo nihil est válidum, nihil san\-ctum:~† multíplica super nos misericórdiam tuam; ut, te re\-ctóre, te duce, sic transeámus per bona temporália,~* ut non amittámus ætérna. Per Dóminum.}{english}
{\noindent
O God, the protector of all them that trust in thee, without whom nothing is strong, nothing is holy; increase and multiply upon us thy mercy, that thou being our ruler and guide, we may so pass through things temporal, that we finally lose not the things eternal. Through our Lord.}

\rubrique{On Saturday, a commemoration of the octave is made: Ant. \normaltext{Ad Jesum, \pageref{ad_jesum_autem}, ℣. Hauriétis Aqua \pageref{haurietis},} and Collect \normaltext{Deus, qui nobis in corde, \pageref{oratio_sc}.}}

    \bigtitle{Sunday within the octave of the Sacred Heart.}
        {\centering which is the III after Pentecost.\par}
\addcontentsline{toc}{section}{Sunday within the octave of the Sacred Heart} \phantomsection

\rubrique{Chapter \normaltext{Humiliámini,} as on Saturday, \normaltext{\pageref{cap_dom_infra_oct_sc}.}}

\rubrique{Hymn \normaltext{En ut supérba críminum, \pageref{hy_en_ut_superba_solesmes_1961}.}}

\rubrique{℣. \normaltext{Memóriam, \pageref{versiculus_dom_octave_sc}.}}

          \gscore[Ad Magnif.]{Ant. 6. f}{an_quae_mulier_solesmes_1961}
        \begin{enpars}\noindent Ant. What woman having ten pieces of silver, if she lose one piece, doth not light a candle, and sweep the house, and seek diligently till she find it?
\end{enpars}

  \rubrique{Collect \normaltext{Proté\-ctor in te sperántium,} as above, \normaltext{\pageref{versiculus_dom_octave_sc}.}}
  
     \rubrique{\begin{enpars} A commemoratio of the following day within the octave is made: Ant. \normaltext{Ignem, \pageref{Ignem}, ℣. Tóllite \pageref{tollite},} and Collect \normaltext{Deus, qui nobis in corde, \pageref{oratio_sc}.}\end{enpars}}

     \rubrique{But if the following day the office to be said is not that of the Octave: Ant. \normaltext{Ad Jesum, \pageref{ad_jesum_autem}, ℣. Hauriétis Aqua \pageref{haurietis},} and collect \normaltext{Deus, qui nobis in corde, \pageref{oratio_sc}.}}
     
     \rubrique{\begin{enpars}The office is as on the feast unless a feast of IX lessons is celebrated instead, and the days within the octave are of semidouble rite.\end{enpars}}

\rubrique{\begin{enpars}On Thursday, the office is of double major rite; all is as at I Vespers of the feast. On Friday, all is as at II Vespers.\end{enpars}}

\end{document}